% commentaire.tex — f167v, vv. 423–465 (conversion initiale des travaux antérieurs)
\comm{424}{ὑποῖξαι τὴν θύραν.}{Il s'agit de la porte de l'appartement (cf. v. 414) et non de celle du cellier.}
\comm{426}{ᾡκότριψ Εὐριπίδης.}{Ce terme, \italique{tâte-poule}, dont un Grec aurait saisi le sens métaphorique (cf. les verbes βλιμάζειν, σκιμαλίζειν), rend bien le mot ᾡκότριψ. Cf. ἀστύτριψ, \italique{pilier de cité}, Θησειότριψ, \italique{pilier du Théseion} (fragm. du Polyides), πορνότριψ, etc. (Note de Willems.)}
\comm{429}{ὄλεθρόν τιν᾽ ἡμᾶς κυρκανᾶν ἁμωσγέπως.}{Κυρκανᾶν est un synonyme très rare de κυκᾶν et se dit proprement d'un cuisinier préparant une mixture\nc{Cf. Épinicus 1,8 : συγκυρκανᾶν (préparation d'une boisson faite de farine, de miel et de vin mêlés).}. Il est pris ici au sens de \italique{comploter la perte} de quelqu'un. Il est clair que l'expression ὄλεθρόν κυρκανᾶν est faite sur le modèle de κυρκανᾶν φάρμακα (cf. \italique{Ploutos} 302, 309 : τὰ φάρμακ᾽ ἀνακυκῶσαν) [...] Pour le ton, le verbe κυρκανᾶν paraît correspondre à nos métaphores familières \italique{cuisiner}, \italique{mitonner une affaire}. (Note de Taillardat, §421.)}
\comm{432}{τῆς γραμματέως.}{Aristophane faisait-il sourire sa mâle assemblée lorsqu'il parlait de \italique{la} greffièr\italique{e} ? c'est ce que suggère Pollux (IV, 19) :\citComm{Ἀριστοφάνης δὲ παίζει ἐν Θεσμοφοριαζούσαις, λέγων ἡ γραμματεύς.}}
\comm{433–42}{}{Ce chant du chœur trouvera son écho aux vers 520-530 : semblable réplique au point de vue métrique\nc{Pour le texte du v. 437b, voir le commentaire critique.} qui permettait aux deux passages d'être chantés sur la même mélodie. Là comme ici, le chœur s'exprime à la première personne.}
\comm{441}{Ξενοκλῆς ὁ Καρκίνου.}{Xénoclès, le plus célèbre des trois carcinites était un poète tragique non dépourvu de talent puisqu'en 415 il l'emporta sur Euripide\nc{Aristophane se montre toutefois sévère avec Xénoclès : il est méchant comme ses vers (v. 169) et son art comme sa personne ne méritent pas de survivre (\italique{Grenouilles}, v. 86).}. Mais on lui reprochait de chercher à captiver son auditoire par des moyens peu relevés. Platon le comique l'appelle δωδεκαμήχανος [l'homme] \italique{aux douze postures}, surnom qu'on donnait à la courtisane Cyrène\nc{Voir la note au vers 235.}. Aristophane le traite de μηχανοδίφης, inventeur de trucs. Il résulte d'un passage d'Athénée (IV, 134d), mis en rapport avec le nôtre, que Xénoclès cumulait avec la profession d'acteur et de poète tragique, celle d'orateur. Ce fut aussi le cas d'Eschine. Cela n'a rien d'étonnant, le métier d'acteur étant une excellente préparation à celui d'orateur. Il n'y a qu'à se souvenir de Démosthène et de ses cailloux. (Note de Willems ; voir aussi la note (t. I pp. 553-558) qu'il a consacrée au banquet de Philocléon conté par Xanthias (\italique{Guêpes}, vv. 1299-1323) où il est question des carcinites.)}
\comm{442}{πᾶσιν ὑμῖν.}{Les spectateurs.}
\comm{446}{ἐμοὶ γὰρ ἁνὴρ ἀπέθανεν μὲν ἔγκυος.}{Pour le texte, voir le commentaire critique.}
\comm{450–1}{ἐν ταῖσιν τραγῳδίαις ποιῶν / τοὺς ἄνδρας ἀναπέπεικεν οὐκ εἶναι θεούς·}{Petit florilège des \italique{\grec{γνῶμαι}} théoclastes d'Euripide :\\
\italique{Hécube}, v. 489-492 :\citComm{Ταλθύβιος — ὦ Ζεῦ, τί λέξω; πότερά σ᾽ ἀνθρώπους ὁρᾶν;\\
ἢ δόξαν ἄλλως τήνδε κεκτῆσθαι μάτην,\\
ψευδῆ, δοκοῦντας δαιμόνων εἶναι γένος\\
τύχην δὲ πάντα τἀν βροτοῖς ἐπισκοπεῖν;}\italique{Bellérophon}\nc{Deux premières lignes du \nieme{10} fragment de l'édition Jouan-van Looy aux Belles Lettres ; fr. 286 de l'édition Kannicht ; fr. 385 Mette.} :\citComm{Βελλεροφόντης — φησίν τις εἶναι δῆτ᾽ ἐν οὐρανῷ θεούς;\\
οὐκ εἰσίν, οὐκ εἴσ᾽,}Prologue de \italique{Mélanippe la sage}\nc{Fragment douteux car on recense différentes versions de ce début de prologue. Citation apocryphe ? retouche d'Euripide d'un vers ayant fait scandale ? quoi qu'il en soit, Jouan \& van Looy le mettent en tête de leur \italique{Mélanippe} sans le comptabiliser ; fr. 480 de l'édition Kannicht.} (douteux) :\citComm{Ζεὺς ὅστις ὁ Ζεύς, οὐ γὰρ οἶδα πλὴν λόγῳ\\
La véritable version euripidéenne étant peut-être\nc{Premier fragment de l'édition Jouan-van Looy ; 481 chez Kannicht et 665a-c chez Mette.} :\\
Μελανίππη — Ζεύς, ὡς λέλεκται τῆς ἀληθείας ὕπο,\\
Ἕλλην᾽ ἔτιχθ᾽, ὃς ἐξέφυσεν Αἴολον·}}
\comm{455–6}{Dossier sur la mère d'Euripide, la vendeuse d'herbes (λαχανοπωλήτρια).}{(\italique{Acharn.} vv. 478-479 ; \italique{Cavaliers} vv. 19-20 ; \italique{Thesm.} vv. 455-456, 910 ; \italique{Grenouilles} vv. 840, 946-947.)


Aristophane plaisante Euripide à maintes reprises sur la profession de sa mère, vendeuse d'herbes\nc{\italique{Thesm.} vv. 387, 456 ; \italique{Acharn.} v. 478 ; \italique{Grenouilles} v. 840.}. Il fit ses choux gras d'une herbe en particulier : ἡ σκάνδιξ, \italique{le scandix peigne de Vénus}. Les commentateurs, anciens comme modernes, se sont interrogés depuis lors pour savoir s'il fallait prendre littéralement cette information, ou bien de façon imagée.
Sans parvenir à aucune conclusion convaincante ; la raison en est que, pour comprendre les attaques d'Aristophane, il fallait séparer le problème en deux :\\
--- quand il évoque le σκάνδιξ, Aristophane moque le style d'Euripide, tout en finasseries (nous allons voir que le verbe διασκανδικίζειν a le même double sens que \italique{épicer} en français ; le sens propre : on épice une galette ; le sens imagé : on épice un discours) ;\\
--- quand il évoque la mère d'Euripide, la verdurière, Aristophane fait entendre que ce dernier est un \italique{fils de la terre}, \latin{terrae filius} ; dont les origines sont obscures et l'arbre généalogique impossible à établir ; bref : Euripide est un bâtard !


Cette note risque fort de trop tirer en longueur ; le dossier sur la mère d'Euripide et son σκάνδιξ est en effet assez épais. En voici le plan :\\
--- je montre d'abord que le vers 910 des \italique{Thesmophories} ne fait pas partie de ce dossier et que les `ἰφύων', \italique{lavandes}\nc{Dans un article passionnant sur quelques plantes apparaissant chez Aristophane, mêlant avec un rare bonheur de savantes études de textes et des connaissances botaniques très pointues, la regrettée Suzanne Amigues identifiait en fait l'ἴφυον avec la nielle des blés (\latin{Agrostemma githago L.}). (« Quelques légumes de disette chez Aristophane et Plutarque » paru en 1988 dans le \italique{Journal des Savants}, pp. 157-171.)}, n'ont jamais apparu dans le texte d'Aristophane ;\\
--- il y a par contre une autre herbe que le scandix à rétablir au vers 20 des \italique{Cavaliers} : l'ἀπόκυνον (Hésychius α 6416) ;\\
--- le σκάνδιξ et le verbe διασκανδικίζειν renvoient au style et à la façon de parler ;\\
--- descendre d'une verdurière, c'est être un \italique{fils de la terre}, \latin{terrae filius} : un bâtard ;\\
--- c'est la bâtardise d'Euripide qui explique sa rancœur vis-à-vis des femmes ; cela ressort de \italique{Thesm.} vv. 455-456 grâce à l'identification d'un proverbe (αὐτὸς τραφείς...) analogue à notre \italique{tel père, tel fils} ;\\
--- la note se conclue par une tentative pseudétymologique sur le σκάνδιξ pour expliquer les images qu'il a générées.



Pour le vers 910 \italique{Thesmophories}, je renvoie le lecteur au commentaire \italique{ad locum}. 
Il n'y a aucune allusion aux herbes d'Euripide, qui ont envahi bien mal à propos jusqu'aux sources les plus fiables de ce vers.

Par contre, il faut en rétablir une aux vers 19-20 des \italique{Cavaliers}. La pièce débute par la traditionnelle scène des deux esclaves -- les généraux Nicias et Démosthène en l'occurrence, généraux qui avaient préparé la victoire de Sphactérie (victoire éclatante par la capture de la garnison spartiate). 
Mais ces deux généraux s'étaient fait voler la vedette par le braillard Cléon débarqué avec un contingent sur place au dernier moment et qui – coup du sort – l'avait emporté\nc{Pour plus de détails, le lecteur pourra utilement se reporter au quart livre de Thucydide, chapitres 28 et suivants ainsi qu'à la longue tirade du Démosthène des \italique{Cavaliers}, à partir du vers 40.}. Ils ressassent leur amertume au début de la pièce en se renvoyant la balle pour trouver une idée de salut : que faire ? — Dis-le ! — Non, toi dis-le\nc{Aristophane fait-il allusion à un complot qui aurait été tramé contre Cléon ? contre ce va-t-en-guerre de Cléon qui empêchait toute négociation de paix au moment-même où Athènes était en position de force pour se sortir favorablement de cette guerre fratricide qui n'avait que trop duré.}... Nicias s'interroge : comment tourner cela précieusement à la manière d'Euripide (κομψευριπικῶς). Démosthène rétorque alors (vers 19-20) :
\citComm{Δημοσθένης --- μὴ 'μοί γε, μὴ 'μοί, μὴ διασκανδικίσῃς·\\
ἀλλ᾽ εὑρέ τιν᾽ †ἀπόκινον† ἀπὸ τοῦ δεσπότου.\\
Démosthène --- Non pas, non pas, ne saupoudre pas de scandix !\\
mais va trouver quelque .......  qui le décroche du maître.}Le « maître », c'est Δῆμος Πυκνίτης, le Populo-citoyen-roi-de-l'Assemblée, mené par le bout du nez par son « chien » : l'infâme Cléon. 
(Tout au long des \italique{Cavaliers}, Cléon est renvoyé à son canin statut de \italique{Chien de Dêmos}, animal flatteur et cajoleur – sûrement par reprise ironique de quelque propos du véritable Cléon, qui se prenait pour le chien de garde de la cité. Voir aussi les \italique{Guêpes}, vv. 852-853, avec le procès aux deux chiens.)


Que le lecteur m'accorde pour l'instant que le verbe διασκανδικίζω signifie « assaisonner à l'aide de σκάνδιξ » et, dans un sens imagé, « agrémenter », « enjoliver [un discours de circonlocutions] » ; le σκάνδιξ jouant un rôle analogue à notre \italique{piment}, et les multiples sens du verbe « pimenter » (Je le vais montrer tout-à-l'heure.)

L'ἀπόκινος des manuscrits, défini par le scholiaste comme une danse licencieuse et qui fait entendre par son étymologie le verbe ἀπο-κινέω, \italique{se retirer}, n'a rien à faire ici\nc{Il faut donc rejeter la jolie trouvaille du grand Alphonse Willems : le « branle de sortie »...}.

Comme toujours devant les difficultés, nous savons que nous pouvons compter sur notre bon vieil et fidèle ami : Hésychius. Voici ce que ce grand érudit a répondu à nos questions (α 6416) :\citComm{ἀπόκυνον: μάζα μεμαγμένη φαρμάκῳ, πρὸς ἀναίρεσιν κυνῶν. ἢ εἶδος βοτάνης\\
apokynon : galette préparée avec du poison [$=$ l'herbe dont il est question juste après], pour la destruction des chiens. \italique{Item} : sorte d'herbe.}


Hésychius nous explique que l'ἀπόκυνον, comme le σκάνδιξ, est une \italique{herbe sauvage}... à cette différence que celui-ci sert de condiment aux pauvres\nc{Le scandix semble être une épice peu appréciée : Andocide parle ainsi des privations des premières années de la guerre du Péloponnèse : μηδὲ ἄγρια λάχανα καὶ σκάνδικας ἔτι φάγοιμεν, « puissions-nous ne plus manger des herbes sauvages ni de scandix ».} quand celui-là nous débarrasse des nuisibles : chiens selon Hésychius, Cléon pour Aristophane. 
Je propose donc la lecture suivante pour les vers 19-20 des \italique{Cavaliers} :
\citComm{Δημοσθένης --- μὴ 'μοί γε, μὴ 'μοί, μὴ διασκανδικίσῃς·\\
ἀλλ᾽ εὑρέ τιν᾽ ἀπόκυνον ἀπὸ τοῦ δεσπότου.\\
(ἀπόκινον codd. omn. : \lecon{ἀπόκυνον} Bibi ap. Hesych.)\\
Démosthène --- Arrête-moi ça : ne condiments pas !\\
mais trouve-moi plutôt quelque mort-aux-chiens qui en débarrasse son maître.}Ce passage des \italique{Cavaliers} est intéressant pour la bonne interprétation du verbe διασκανδικίζειν : Démosthène l'utilise d'abord comme synonyme de parler précieusement à la manière d'Euripide (κομψευριπικῶς) ; on dirait en français « épicer un discours ».  Mais au vers 20, il le sous-entend avec son sens premier : « Va me trouver de l'ἀπόκυνον [pour m'en épicer une galette] qui nous débarrassera le maître [de ce chien de Cléon]. »

Le verbe διασκανδικίζω n'est pas un hapax d'Aristophane ; il se trouve aussi dans un fragment de Téléclide, un autre poète comique contemporain d'Euripide. Ce fragment\nc{Fragment 40(38) dans l'édition de Kassel-Austin, \italique{Poetae Comici Graeci}, vol. VII.} nous a été conservé par Athénée (\italique{Deipnosophistes}, VI, 56d.) :
\citComm{Ξυγγενέσθαι διὰ χρόνου †λιπαρείτω† με\\
δρυπεπέσι μάζαις καὶ διασκανδικίσαι.}
Tout ce que l'on peut inférer du contexte chez Athénée, c'est que l'adjectif δρυπεπής, \italique{bien mûr}, lorsqu'il est substantivé comme ici, désigne une \italique{olive}. 
Le sens que j'ai donné à διασκανδικίζω, « assaisonner à l'aide de σκάνδιξ », rendrait bien mieux compte aussi de ce second passage que celui indiqué dans les dictionnaires : « se nourrir de σκάνδιξ »\nc{« manger du cerfeuil » Bailly ; « \espagnol{dar de comer perifollo} » DGE ; « \anglais{feed on chervil} » LSJ ; « Kerbel essen » Pape.}.
Je traduis donc : « Il m'importune (λιπαρεῖ ?) depuis longtemps pour ajouter des olives aux galettes ainsi que les agrémenter de scandix ».

Je n'ai pas encore parlé d'une citation non identifiée à l'entrée Σκάνδιξ de la \italique{Souda} (σ 536) qui montre clairement que le scandix, en un sens imagé, désigne un \italique{euphémisme} :
\citComm{καὶ μή μοι τὰ τοῦ σκανδικοπώλου παρατραγῴδει, σεμνὸν τὸ τῆς Πειρήνης ὕδωρ ἀποκαλοῦντος,\\
« et ne m'use pas de l'emphase du vendeur de scandix qui appelle `sacrée' la fontaine de Pirène ».}La cible est clairement Euripide : la citation peut faire allusion au vers 69 de sa \italique{Médée} ou aux vers 205-206 des \italique{Troyennes}. Il est toujours délicat de proposer des hypothèses sur des fragments aussi brefs... proposons tout-de-même notre hypothèse qui vaut ce qu'elle vaut – c'est-à-dire pas grand-chose : cette fontaine située sur l'Acrocorinthe délivrait effectivement une eau pure et abondante\nc{Cf. ce qu'en disent Pausanias, Strabon et Suidas à l'entrée Πειρήνη ainsi que cette citation d'Antipater prise dans l'\italique{Anthologie grecque} : Λαί̈δ' ἔχω πολιῆτιν ἐϋζώνοιο Κορίνθου, Πειρήνης λευκῶν φαιδροτέρην λιβάδων, « Je tiens Laïs, citoyenne de Corinthe, bien ceinturée, plus brillante que les gouttes blanches de Pirène ».} ; ce n'est donc pas sur la qualité de l'eau qu'Euripide enrobe de scandix la vérité... peut-être sur ce qu'on y pratiquait ? La fontaine de Pirène se trouvait juste en contrebas du temple d'Aphrodite dont certaines de ses préposées se livraient à une prostitution « sacrée »\nc{La réputation des Corinthiennes était telle que les Grecs en ont inventé le verbe κορινθιάζειν.}. Imagine-t-on un tragique contemporain, si peu tragique qu'il nous parlerait du « \italique{sacré} bois de... Boulogne » (!?) ce serait un peu fort de scandix !.. Cet exemple nous dessine très clairement le sens imagé du scandix : Euripide \italique{dore la pilule} en saupoudrant de scandix des choses basses ; il nous fait passer des vessies pour des lanternes.

Voici pour le style d'Euripide. Passons à sa généalogie. L'on connaît le \italique{topos} comique : chaque fois qu'un poète voulait discréditer l'un de ses contemporains, il l'accusait de n'être pas de pure race (εὐγενής) : soit qu'il fut Athénien d'importation récente\nc{Voir par exemple \italique{Grenouilles} vv. 717-733 qui est sans ambiguïté. L'on s'étonne de trouver dans les modernes traductions des \italique{Thesmophories} (v. 330) εὐγενεῖς γυναῖκες rendu par `nobles femmes', comme si la possibilité de célébrer les Thesmophores était conditionnée à une forme quelconque de noblesse. Il serait bon avant d'aborder Aristophane de se défaire de nos préjugés modernes, faute de quoi l'on s'expose à de tels contre-sens. Lire là-dessus l'indispensable article d'Alphonse Willems « Aristophane et la démocratie athénienne » reproduit en annexe et qu'on peut trouver dans son \italique{Aristophane} posthume (1919), t. III, pp. 445-473.}, soit qu'on puisse émettre des doutes sur sa lignée officielle. Euripide n'y échappe pas. Dans les \italique{Grenouilles} (vv. 946-947), alors qu'il expose sa théorie dramatique, il est interrompu :\citComm{Εὐριπίδης --- εἶτ’ οὐκ ἐλήρουν ὅ τι τύχοιμ’ οὐδ’ ἐμπεσὼν ἔφυρον,\\
ἀλλ’ οὑξιὼν πρώτιστα μέν μοι τὸ γένος εἶπ’ ἂν εὐθὺς\\
τοῦ δράματος.\\
Διόνυσος --- \decalage{τοῦ }κρεῖττον γὰρ ἦν σοι νὴ Δί’ ἢ τὸ σαυτοῦ.\\
Euripide --- Je ne débitais pas toutes les niaiseries qui me venaient, et je ne brouillais pas tout, sans crier gare ;
mais le premier qui entrait en scène exposait de prime-saut la filiation du drame.\\
Dionysos --- Cela valait mieux pour toi, par Zeus, que d'exposer la tienne.}L'attaque n'est pas neuve : quelques vers plus tôt, Euripide avait provoqué Eschyle ; la réplique ne s'était pas fait attendre (vv. 836-843) :\citComm{Εὐριπίδης --- ἐγᾦδα τοῦτον καὶ διέσκεμμαι πάλαι,\\
ἄνθρωπον ἀγριοποιὸν αὐθαδόστομον,\\
ἔχοντ᾽ ἀχάλινον ἀκρατὲς ἀπύλωτον στόμα,\\
ἀπεριλάλητον κομποφακελορρήμονα.\\
Αἰσχύλος --- `ἄληθες ὦ παῖ τῆς ἀρουραίας θεοῦ;'\\
σὺ δή με ταῦτ᾽ ὦ στωμυλιοσυλλεκτάδη\\
καὶ πτωχοποιὲ καὶ ῥακιοσυρραπτάδη;\\
ἀλλ᾽ οὔ τι χαίρων αὔτ᾽ ἐρεῖς.\\
Euripide --- Je le connais et l'ai de longtemps percé à jour,\\
cet homme aux héros sauvages, au verbe présomptueux,\\
à la langue sans frein, sans retenue, sans barrières,\\
à l'esprit rétif, ce fagoteur de mots ronflants.\\
Eschyle --- « Ah, par exemple ! fils de la déesse agreste », \\
c'est ainsi que tu parles de moi, raccoleur de fadaises,\\
faiseur de mendiants, ravaudeur de haillons.\\
Mais tu n'auras pas lieu de t'en applaudir.}
(Le scholiaste nous indique que la réaction cinglante d'Eschyle au vers 840 est une parodie d'un vers d'Euripide avec ἀρουραίας, \italique{agreste}, en place de θαλασσίας, \italique{marine}.)



Or on désigne comme fils de la (déesse) terre dans la langue proverbiale grecque (et latine) les personnes dont la généalogie est obscure. 
Dans la pièce d'Euripide qui porte son nom, Ion interroge Xouthos sur l'identité de sa mère ; comme Xouthos ignore la réponse, Ion s'exclame alors (v. 542) :\citComm{Ἴων --- γῆς ἄρ᾽ ἐκπέφυκα μητρός;\\
Ion --- J'ai donc eu la terre pour mère ?}
Érasme a recensé l'adage (n° 786) latin \latin{terrae filius} employé avec le même sens par Tertullien, Cicéron, \&c.
C'est bien le propre de l'αὐτό-χθων que de ne pas savoir à quand remonte l'installation de ses ancêtres dans le pays : sa filiation se perd dans la nuit des temps.

L'on comprend maintenant où veut en venir Aristophane avec ses légumes et autres herbes, étroitement associés à la terre : si nul Français ne s'offusquera qu'on dise de lui qu'il est né dans un chou ou une rose, en grec cela revient à émettre des doutes sur la naissance.
Ainsi s'explique notre passage des \italique{Thesmophories} (vv. 455-456).
Le procès d'Euripide a été ouvert et la parole est à l'accusation. Après un premier réquisitoire-fleuve, une seconde oratrice s'avance et prend la parole (vers 443-458).
Le point culminant de l'attaque se trouve dans le distique des vers 455-456 :\citComm{ἄγρια γὰρ ἡμᾶς ὦ γυναῖκες δρᾷ κακά,\\
ἅτ᾽ ἐν ἀγρίοισι τοῖς λαχάνοις αὐτὸς τραφείς.}Aristophane joue sur la répétition d'ἄγριος, qualifiant les mauvais comportements qu'Euripide impute à l'engeance féminine dans ses pièces\nc{Aristophane, au vers 349 des \italique{Nuées}, qualifie ainsi d' ἄγριος le chevelu Hiéronyme, un poète dithyrambique, fils de Xénophane qui passait pour pédéraste ; et le scholiaste de noter qu'on appelait ἀγρίους (et débauchés) les pédérastes : ἀγρίους (καὶ κολλοποδιώκτας) ἐκάλουν τοὺς παιδεραστάς. De même, à l'entrée Ἀγρίους de la \italique{Souda}, on peut lire : Αἰσχίνης ἐν τῷ κατὰ Τιμάρχου τοὺς σφόδρα ἐπτοημένους περὶ τὰ παιδικὰ οὕτως ὠνόμασε. καὶ Μένανδρος δὲ ἄγριον εἶπε κυβευτὴν τὸν λίαν περὶ τὸ κυβεύειν ἐσπουδακότα, « Eschine dans son \italique{Contre Timarque} nomme ainsi les hommes en proie à une excessive attirance pour les enfants. Et Ménandre dit qu'un joueur est ἄγριον, celui qui est trop porté au jeu ». L'on voit que le terme ἄγριος, au sens moral, désigne un \italique{intempérant}.} et, au vers suivant, les λάχανα, les \italique{légumes} qui, lorsqu'ils sont ἄγρια, « sauvages », désignent des \italique{herbes sauvages}, de celles qui poussent sans qu'on les cultive (par opposition à ἥμερα). 

Si Euripide nourrit une telle haine vis-à-vis des femmes et les présente comme dépravées, c'est que sa propre mère était telle... Sur le fond, ce genre de psychanalyse de comptoir n'est pas d'une grande finesse, mais le coup sous la ceinture porte auprès des femmes et suffit pour déchaîner l'enthousiasme du chœur quelques vers plus loin\nc{κομψότερον ἔτ᾽ ἢ τὸ πρότερον, « plus intelligent encore que le premier [discours] » [...] οἷα κατεστωμύλατο, « comme elle a bien parlé » [...] φρένας ἔχουσα, « elle a de l'esprit » \&c.}.

Ces éloges concernent davantage la forme que le fond de la diatribe\nc{L'oratrice précédente, au vers 387, avait déjà traité Euripide de fils de la λαχανοπωλήτρια. Plutarque citait ces deux vers comme exemple d'un jeu de mots qu'il goûtait fort peu (voir ses \italique{Œuvres morales}, 853C, \italique{Abrégé de la comparaison d'Aristophane et de Ménandre}). Je ne suis pas sûr que cet inénarrable raseur, le Cicéron des lettres grecques, ait bien saisi intégralement le détournement d'Aristophane.}.

Le tour `αὐτὸς τραφείς' est une locution proverbiale qu'on retrouve ailleurs : le parémiographe Michel Apostolius\nc{\italique{Corpus paroemiographorum graecorum}, II, ed. Leutsch, Apostolii centuria V 30a, p. 341.} a relevé dans l'\italique{Anthologie} de Stobée (IV, 25, 6), au chapitre consacré au respect filial, cette citation du poète Dikaiogène\nc{L'on ne possède pas grand-chose sur ce poète grec, auteur de tragédies et dithyrambes, qui vécut probablement au \nieme{IV} siècle av. J.-C. Tout ce que l'on sait de lui est regroupé dans les \italique{Tragicorum Graecorum Fragmenta}, vol. I, ed. Snell \& Kannicht. Il y est classé en tant que « poète mineur » n° 52.} :\citComm{Αὐτὸς τραφεὶς δὲ τῶν φυτευσάντων ὕπο\\
καλῶς τὸν αὐτὸν ἔρανον\nc{L'on trouve chez Démosthène dans son \italique{Contre Neæra}, 8, l'expression ironique τὸν αὐτὸν ἔρανον ἀποδοῦναι, \italique{payer quelqu'un de la même monnaie}.} αὐτοῖσιν νέμεις\\
Toi-même élevé [litt. : nourri] par tes parents [litt. : ceux qui t'ont semé],\\
tu leur partages comme de juste ton propre repas.}
La locution proverbiale complète était-elle `αὐτὸς τραφεὶς ὑπὸ τῶν φυτευσάντων' ? On comprendrait où veut en venir l'oratrice : si les parents d'Euripide, ses \italique{semeurs}, l'ont semé au milieu des légumes, il n'est pas douteux que cet homme de la terre y ait récolté tout le fiel qu'il déverse sur les femmes dans ses tragédies.

Dion Cassius (\italique{Histoire romaine}, XLVI, 4) met dans la bouche d'un orateur la même figure de rhétorique, semblable par la forme (ἐν $+$ datif $+$ αὐτὸς τραφείς) et par le ton. C'est une réponse à Cicéron et ses violentes \italique{Philippiques} contre Marc Antoine. Cicéron... qui est-ce, au juste ?
\citComm{ὁ πατὴρ αὐτῷ ὁ κναφεύς [ ... ] ἄνθρωπος ἀγαπητῶς ἔκ τε τούτων καὶ ἐκ τῶν πλυνῶν διατρεφόμενος [ ... ] \textbf{Ἐν οἷς αὐτὸς τραφεὶς} οὐκ ἀπεικότως τοὺς κρείσσονας αὐτοῦ καὶ πατεῖ καὶ πλύνει.\\
\francais{Son père était foulon, lui qui avait grand-peine à vivre de ce métier et de son lavoir. À la suite d'une telle éducation, non sans invraisemblance, ce personnage foule aux pieds ceux qui valent mieux que lui et leur lave la tête.}}
Ces clefs d'interprétation permettent d'expliquer le vers 478 des \italique{Acharniens}, alors que Dikaiopolis emprunte à Euripide tout un fatras de matériel puisé dans ses propres drames, voilà qu'il déclenche le courroux d'Euripide en lui demandant :
\citComm{Δικαιόπολις --- σκάνδικά μοι δὸς μητρόθεν δεδεγμένος.\\
Εὐριπίδης --- ἁνὴρ ὑβρίζει· κλῇε πηκτὰ δωμάτων.\\
Dikaiopolis --- Donne-moi de ce scandix que tu tiens de ta mère.\\
Euripide --- Il m'insulte. « Ferme sur lui les ais solides de la porte. »}
Aristophane fait coup double : l'on comprend bien pourquoi Dikaiopolis veut emprunter à Euripide son scandix : il a besoin de la faconde du tragédien pour la joute verbale qui va l'opposer au chœur des vieillards d'Acharnes\nc{C'est la raison de sa visite chez Euripide (vers 416-417) : δεῖ γάρ με λέξαι τῷ χορῷ ῥῆσιν μακράν· /
αὕτη δὲ θάνατον, ἢν κακῶς λέξω, φέρει. 
« Car j'ai à débiter au chœur une longue tirade. /
Il m'en coûte la vie, si je m'en tire mal. »}. 
On pourrait dire que Dikaiopolis compte les « enfariner » (διασκανδικίζειν) avec l'art oratoire d'Euripide... mais en précisant que ce scandix lui vient de sa mère, Dikaiopolis fait allusion au sobriquet de fils de verdurière qui colle au poète.
Euripide le comprend bien qui chasse l'insolent : ἁνὴρ ὑβρίζει· κλῇε πηκτὰ δωμάτων.\\

Cette trop longue note touchant enfin à sa fin, je me permettrai quelque aventureuse conjecture sur l'origine de cette utilisation imagée du σκάνδιξ : il me semble que si le grec a choisi cette plante précisément et non une autre, c'est non pas pour ses vertus (réelles ou supposées), mais bien plutôt pour des raisons de sonorité. Cette intuition m'est venue de la consultation de l'\latin{Etymologicum Genuinum}\nc{Lexique du \nieme{IX} siècle souvent associé à d'autres, postérieurs mais plus complets : l'\latin{Etymologicum magnum} et l'\latin{Etymologicum Symeonis}.} à l'entrée σκάνδαλον, la \italique{chausse-trappe}, pris ici en un sens imagé : σκάνδαλον: τὸ σκάζειν περὶ τὴν ἀλήθειαν. 
Image saisissante et évocatrice s'il en est ! « boîter autour de la vérité », « tourner autour du pot-aux-roses ». 
Elle s'appuie sur un rapprochement pseudo-étymologique\nc{Je me vois forcé de préciser \italique{pseudo}-étymologique, car mes velléités de rapprocher étymologiquement σκάνδαλον de σκάζειν furent foudroyées par Alexandre Perret, dont j'eus la chance d'assister à des cours de linguistique, cours qui ressemblaient davantage à d'enivrants voyages à travers les contrées et les parlers d'Europe (et ce jusqu'aux confins de l'Asie). Voici ce que m'a très aimablement répondu M. Perret : « Après recherches fructueuses, je pense avoir trouvé une réponse dans le \italique{Lexikon der Indogermanischen Verben} (LIV) mais il ne semble pas que les deux racines soient liées. Σκανδάλον semble en effet rattaché à un thème *skend- qui signifierait « bondir » (on a skándati en sanskrit pour « sauter »), mais le suffixe -αλον me dérange dans la mesure où ça semble pré-grec. Aucune étymologie n'est assez bien solide pour ce mot qui peut être un emprunt. Quelle motivation pour passer de « bondir » à « faire tomber » ? À la limite, un autre thème *sk'end- (avec un k- palatal donc invisible en grec mais pas en sanskrit) avec le sens parfois de « faire tomber » (skr. Chándayati: « tomber »), est en soit beaucoup plus plausible !! Mais pour σκάζω, on rapproche ça d'une racine *skeng- avec effectivement le sens de « boiteux, tors » (All. \italique{Hinken}). Ce qui m'étonne est qu'Hesychius recense à la fois σκάνδαλον: \italique{piège à souris} et σκάνδαλος: \italique{obstacle}, mais est-ce que les pièges à souris étaient des trappes où tomber, ou bien des claquoirs comme à notre époque ? Ça mérite investigation. En soit passer de *skend- à σκάζω avec un suffixe -yo de présent n'est pas choquant. »} entre le mot σκάνδ-αλον et le verbe σκάζειν, boîter... le σκάνδ-ιξ ne pouvait-il pas suggérer la même image : celle de \italique{tordre}, \italique{rendre boiteuse} ou de \italique{camoufler} la vérité en la saupoudrant d'une herbe de vil prix ?


Voilà qui donne quelques nouveaux éléments d'appréciation sur le problème d'Euripide, son σκάνδιξ et sa maman la λαχανοπωλήτρια.}
\comm{458}{στεφάνους συνθηματιαίους.}{Athénée (XV, 680c), dans le long chapitre consacré aux couronnes de ses \italique{Deipnosophistes}, parle de ces couronnes \italique{synthématiées} :\citComm{συνθηματιαῖοι στέφανοι, ἠργολαβημένοι καὶ ἐκδόσιμοι. Ἀριστοφάνης ἐν Θεσμοφοριαζούσαις.}Glose identique chez Hésychius (σ 2624) :\citComm{συνθηματιαῖοι: οἱ ἠργολαβημένοι καὶ ἐκδεδομένοι, ἐκδόσιμοι.}}
